% ============ 16 纯性; Gysin 序列 (原书 pp.108-113) ============
\chapter{纯性；Gysin 序列}\label{sec:16}
\index{纯性}
\index{Gysin 序列}
\index{Tate 扭转}
\index{完全交}
\index{光滑对}
\index{上同调纯性}

固定整数 $n > 0$，设 $\Lambda = \mathbb{Z}/n\mathbb{Z}$.对 $n \cdot 1$ 是 $R$ 中单位的任何环 $R$，定义 $\mu_n(R)$ 为 $R$ 中 $1$ 的 $n$ 次单位根群，并定义
\[
\mu_n(R)^{\otimes r} = \begin{cases}
\underbrace{\mu_n(R) \otimes \cdots \otimes \mu_n(R)}_{r \text{ 个}}, & r > 0;\\
\Lambda, & r = 0;\\
\operatorname{Hom}_{\Lambda}(\mu_n(R)^{\otimes (-r)}, \Lambda), & r < 0.
\end{cases}
\]
当 $R$ 是含 $n$ 次单位根的整环时，这些群都是 $\Lambda$ 上的秩 1 自由模，而本原 $n$ 次单位根的选取同时确定它们全体的一组基.

设 $X$ 是特征不整除 $n$ 的域 $k$ 上的簇(或满足 $n\mathcal{O}_X = \mathcal{O}_X$ 的概形).我们定义 $\Lambda(r)$ 为 $X_{\mathrm{et}}$ 上满足
\[
\Gamma(U, \Lambda(r)) = \mu_n(\Gamma(U, \mathcal{O}_U))^{\otimes r}
\]
($U \to X$ 平展且仿射)的层.若 $k$ 含 $n$ 次单位根，则每个层都同构于常值层 $\Lambda$，而 $k$ 中本原 $n$ 次单位根的选取确定对所有 $r$ 的同构 $\Lambda(r) \xrightarrow{\sim} \Lambda$.无论如何，每个层 $\Lambda(r)$ 都局部常值.对 $X_{\mathrm{et}}$ 上被 $n$ 消灭的层 $F$，设
\[
F(r) = F \otimes \Lambda(r), \qquad \text{所有 } r \in \mathbb{Z}.
\]

\subsection*{定理的陈述；推论}

为简单起见，现在假定 $k$ 代数闭.

$k$-簇的光滑对 $(Z, X)$ 是非奇异 $k$-簇 $X$ 连同非奇异子簇 $Z$.我们说 $(Z, X)$ 有余维数 $c$，如果 $Z$ 的每个连通分支\footnote{回忆非奇异簇的每个连通分支都不可约.}在 $X$ 的对应分支中都有余维数 $c$.

\begin{theorem}\label{thm:16.1}
对余维数 $c$ 的 $k$-簇的光滑对 $(Z, X)$ 与 $X$ 上的 $\Lambda$-模局部常值层 $F$，对所有 $r \geq 0$ 有典范同构
\[
H^{r-2c}(Z, F(-c)) \to H^{r}_{Z}(X, F).
\]
\end{theorem}

\begin{corollary}\label{cor:16.2}
在定理的情形中，有同构
\[
H^{r}(X, F) \to H^{r}(U, F), \qquad 0 \leq r < 2c - 1,
\]
与正合序列(Gysin 序列)
\[
0 \to H^{2c-1}(X, F) \to H^{2c-1}(U, F|_{U}) \to \cdots
\]
\[
\to H^{r-2c}(Z, F(-c)) \to H^{r}(X, F) \to H^{r}(U, F) \to \cdots.
\]
\end{corollary}

\begin{proof}
用定理把空间对 $(X, X \setminus Z)$ 的正合序列中的群 $H^{r}_{Z}(X, F)$ 换成群 $H^{r-2c}(Z, F(-c))$.
\end{proof}

\begin{example}\label{ex:16.3}
回忆对 $k$ 上任何非奇异仿射曲线 $U$，$H^{1}(U, \mathbb{G}_{m}) = \operatorname{Pic}(U)$，且对 $r > 1$，$H^{r}(U, \mathbb{G}_{m}) = 0$.因此 $H^{1}(\mathbb{A}^{1}, \mathbb{G}_{m}) = 0$，因为 $k[T]$ 是主理想整环.Kummer 序列
\[
\begin{tikzcd}
0 \arrow[r] & H^{0}(\mathbb{A}^{1}, \mu_n) \arrow[r] \arrow[d, equal] & H^{0}(\mathbb{A}^{1}, \mathbb{G}_{m}) \arrow[r] \arrow[d, equal] & H^{0}(\mathbb{A}^{1}, \mathbb{G}_{m}) \arrow[r] \arrow[d, equal] & H^{1}(\mathbb{A}^{1}, \mu_n) \arrow[r] & 0 \\
& k^{\times} \arrow[r, "n"] & k[T]^{\times} \arrow[r, "n"] & k[T]^{\times} &
\end{tikzcd}
\]
的上同调序列表明对 $r > 0$，$H^{r}(\mathbb{A}^{1}, \mu_n) = 0$.K\"unneth 公式(\S 22)现在蕴含对 $r > 0$，$H^{r}(\mathbb{A}^{m}, \Lambda) = 0$，即 $\mathbb{A}^{m}$ 是``无环的''.因此，对 $(\mathbb{P}^{m-1}, \mathbb{P}^{m})$ 的 Gysin 序列表明
\[
H^{0}(\mathbb{P}^{m}, \Lambda) \cong H^{0}(\mathbb{A}^{m}, \Lambda) \cong \Lambda
\]
\[
H^{1}(\mathbb{P}^{m}, \Lambda) \hookrightarrow H^{1}(\mathbb{A}^{m}, \Lambda) = 0
\]
且
\[
H^{r-2}(\mathbb{P}^{m-1}, \Lambda(-1)) \cong H^{r}(\mathbb{P}^{m}, \Lambda).
\]
归纳论证现在表明
\[
H^{r}(\mathbb{P}^{m}, \Lambda) = \Lambda, 0, \Lambda(-1), 0, \Lambda(-2), \ldots, \Lambda(-m),
\]
\[
r = 0 \quad 1 \quad 2 \quad 3 \quad 4 \quad \ldots \quad 2m
\]
即
\[
H^{r}(\mathbb{P}^{m}, \Lambda) = \begin{cases}
\Lambda(-\frac{r}{2}), & r \text{ 偶}, \leq 2m;\\
0, & \text{其他情形}.
\end{cases}
\]
\end{example}

\begin{example}\label{ex:16.4}
设 $X$ 是 $\mathbb{P}^{m+1}$ 中的非奇异超曲面，即齐次理想 $\mathcal{I}(X) = (f)$ 的 $\mathbb{P}^{m+1}$ 的闭子簇，其中 $f$ 是 $k[T_0, \ldots, T_{m+1}]$ 中的齐次多项式，使得
\[
f, \frac{\partial f}{\partial T_0}, \ldots, \frac{\partial f}{\partial T_{m+1}}
\]
在 $\mathbb{P}^{m+1}$ 中没有公共零点.则 $U \overset{\mathrm{def}}{=} \mathbb{P}^{m+1} \setminus X$ 仿射(AG 6.21)，从而对 $r > m + 1$，$H^{r}(U, \Lambda) = 0$.因此，Gysin 序列为我们提供映射
\[
H^{r}(X, \Lambda) \to H^{r+2}(\mathbb{P}^{m+1}, \Lambda(1)),
\]
它们对 $r > m$ 是同构，对 $r = m$ 是满射.因此对 $r > m$，$H^{r}(X, \Lambda) \cong H^{r}(\mathbb{P}^{m}, \Lambda)$，且 $H^{m}(X, \Lambda) \cong H^{m}(\mathbb{P}^{m}, \Lambda) \oplus H^{m}(X, \Lambda)^{0}$，其中 $H^{m}(X, \Lambda)^{0}$ 是 $H^{m}(X, \Lambda) \to H^{m+2}(\mathbb{P}^{m+1}, \Lambda(1))$ 的核.用 Poincar\'e 对偶(见下一个附记)，得到典范分解
\[
H^{*}(X, \Lambda) \cong H^{*}(\mathbb{P}^{m}, \Lambda) \oplus H^{m}(X, \Lambda)^{0},
\]
其中 $H^{*}(\cdot, \Lambda)$ 表示分次 $\Lambda$-模 $\bigoplus H^{r}(\cdot, \Lambda)$.

更一般地，$\mathbb{P}^{N}$ 中维数 $m$ 的闭子簇 $X$ 称为完全交，如果它的齐次理想由 $N - m$ 个多项式生成.则存在链
\[
\mathbb{P}^{N} \supset X_{N-1} \supset \cdots \supset X_m = X,
\]
其中每个 $X_r = H_r \cap X_{r+1}$(作为概形)，$H_r$ 是 $\mathbb{P}^{N}$ 中的超曲面.我们说 $X$ 是光滑完全交，如果存在所有 $X_r$ 都非奇异的这样的链.对维数 $m$ 的光滑完全交 $X$，用 Gysin 序列与 Poincar\'e 对偶作归纳论证再次表明有典范分解
\[
H^{*}(X, \Lambda) \cong H^{*}(\mathbb{P}^{m}, \Lambda) \oplus H^{m}(X, \Lambda)^{0}.
\]
我们后面将看到 $H^{m}(X, \Lambda)^{0}$ 是秩 $\beta'_{m}$ 的自由 $\Lambda$-模，该秩只依赖于 $m$ 与生成 $\mathcal{I}(X)$ 的多项式的次数，并且事实上 $\beta'_{m}$ 有显式公式.例如，若 $X$ 是次数 $d$，维数 $m$ 的非奇异超曲面，则
\[
\beta'_{m} = \frac{(d-1)^{m+2} + (-1)^{m}(d-1)}{d}.
\]
后面我们将看到，这些结果蕴含：若 $X$ 是有限域 $\mathbb{F}_q$ 上的维数 $m$ 的光滑完全交，即生成 $\mathcal{I}(X)$ 的 $N - m$ 个多项式可以选取系数在 $\mathbb{F}_q$ 中，则
\[
|\#X(\mathbb{F}_q) - \#\mathbb{P}^{m}(\mathbb{F}_q)| \leq \beta'_{m} q^{\frac{m}{2}}.
\]
与 $\#X(\mathbb{F}_q) \sim q^{m}$ 相比，误差项 $\beta'_{m} q^{\frac{m}{2}}$ 非常小！

遗憾的是，光滑完全交是非常特殊的簇.例如，当我们爆破维数 $m > 1$ 的非奇异簇 $X$ 上的点 $P$ 时，$P$ 被换成与 $P$ 处的切空间相联系的射影空间，即 $\mathbb{P}^{m-1}$ 的副本.于是，对新簇 $Y$，
\[
\#Y(\mathbb{F}_q) = \#X(\mathbb{F}_q) + \#\mathbb{P}^{m-1}(\mathbb{F}_q) - 1 = \#X(\mathbb{F}_q) + q^{m-1} + \cdots + q.
\]
若上述不等式对 $X$ 成立，则对 $Y$ 不成立.
\end{example}

\begin{aside}\label{asd:16.5}
后面我们将看到如下形式的 Poincar\'e 对偶定理：对 $k$ 上维数 $m$ 的非奇异簇 $U$，有典范完美配对
\[
H^{r}_{c}(U, \Lambda) \times H^{2m-r}(U, \Lambda(m)) \to H^{2m}_{c}(U, \Lambda(m)) \cong \Lambda.
\]
这里 $H^{r}_{c}$ 是``紧支集上同调群''，当 $X$ 完备时它等于通常的上同调群.
\end{aside}

\begin{example}\label{ex:16.6}
设 $X$ 是射影空间中维数 $m$ 的连通非奇异闭簇.由 Bertini 定理(Hartshorne 1977, III.7.9.1)，存在连通非奇异簇的序列
\[
X = X_m \supset X_{m-1} \supset \cdots \supset X_1,
\]
其中每个 $X_r$ 是前一个簇的超平面截.由 Gysin 序列，得到映射
\[
H^{1}(X_1, \Lambda) \xrightarrow{\text{满}} H^{3}(X_2, \Lambda(1)) \xrightarrow{\sim} \cdots \xrightarrow{\sim} H^{2m-1}(X_m, \Lambda(m-1)).
\]
用 Poincar\'e 对偶，得到单射
\[
H^{1}(X, \Lambda(1)) \to H^{1}(X_1, \Lambda(1)). \qquad (*)
\]
包含 $X_1 \hookrightarrow X$ 在 Picard 群上定义映射
\[
\operatorname{Pic}(X) \to \operatorname{Pic}(X_1).
\]
映射 $(*)$ 可以等同于
\[
\operatorname{Pic}(X)^m \to \operatorname{Pic}(X_1)^m.
\]
代数闭域上阿贝尔簇的挠点在 Zariski 拓扑中稠密，因此这表明由 $X_1 \hookrightarrow X$ 诱导的映射
\[
\operatorname{PicVar}(X) \to \operatorname{Jac}(X_1)
\]
是单射.这个结果在经典中已知：对非奇异簇 $X$ 上的``一般''曲线 $C$，$X$ 的 Picard 簇是 $C$ 的 Jacobi 簇的子簇.
\end{example}

\subsection*{定理的重述}

现在允许 $k$ 是任意域.

设 $Z$ 是簇 $X$ 的闭子簇(或概形 $X$ 的闭子概形)，$i \colon Z \hookrightarrow X$ 是包含映射.如我们在 8.18 中看到的，函子 $i_{*}$ 与 $i^{*}$ 定义了支集在 $Z$ 上的 $X$ 的平展层范畴与 $Z$ 的平展层范畴之间的等价.由于 $i_{*}$ 正合且保持内射对象，$H^{r}(X, i_{*}F) = H^{r}(Z, F)$，有时允许把 $F$ 与 $i_{*}F$ 混淆.

如常，设 $U = X \setminus Z$，用 $j$ 表示包含 $U \hookrightarrow X$.

对 $X$ 上的层 $F$，定义 $F^{\dagger}$ 为 $F$ 的支集在 $Z$ 上的最大子层.于是，对任何平展映射 $\varphi' \colon V \to X$，
\[
F^{\dagger}(V) = \Gamma_{\varphi'^{-1}(Z)}(V, F) = \operatorname{Ker}(F(V) \to F(\varphi'^{-1}(U))).
\]
容易看出这确实定义 $X$ 上的层，且
\[
F^{\dagger} = \operatorname{Ker}(F \to j_{*}j^{*}F).
\]
设 $G$ 是支集在 $Z$ 上的 $X$ 上的层.则任何同态 $\alpha \colon G \to F$ 的像都在 $F^{\dagger}$ 中，从而
\[
\operatorname{Hom}_X(G, F^{\dagger}) \cong \operatorname{Hom}_X(G, F).
\]
现在定义 $i^{\dagger}F$ 为看作 $Z$ 上层的 $F^{\dagger}$，即 $i^{\dagger}F = i^{*}F^{\dagger}$.则 $i^{\dagger}$ 是函子 $\mathrm{Sh}(X_{\mathrm{et}}) \to \mathrm{Sh}(Z_{\mathrm{et}})$，而显示的同构可以解释为典范同构
\[
\operatorname{Hom}_Z(G, i^{\dagger}F) \cong \operatorname{Hom}_X(i_{*}G, F).
\]
这表明 $i^{\dagger}$ 是 $i_{*}$ 的右伴随.因为 $i^{\dagger}$ 有左伴随，它左正合，且因为左伴随正合，它保持内射对象.

\begin{theorem}\label{thm:16.7}[上同调纯性]
设 $(Z, X)$ 是余维数 $c$ 的代数簇的光滑对.对 $X$ 上的 $\Lambda$-模局部常值层 $F$，$R^{2c}i^{\dagger}F \cong (i^{*}F)(-c)$，且对 $r \neq 2c$，$R^{r}i^{\dagger}F = 0$.
\end{theorem}

\textbf{为什么定理~\ref{thm:16.7} 蕴含定理~\ref{thm:16.1}} 考虑函子
\[
\mathrm{Sh}(X_{\mathrm{et}}) \xrightarrow{i^{\dagger}} \mathrm{Sh}(Z_{\mathrm{et}}) \xrightarrow{\Gamma(Z, \cdot)} \mathbf{Ab}.
\]
这些函子左正合，$i^{\dagger}$ 保持内射对象，而复合是 $\Gamma_Z(X, \cdot)$.因此(见定理~\ref{thm:10.7})，存在谱序列
\[
E^{r, s}_2 = H^{r}(Z, R^{s}i^{\dagger}F) \Rightarrow H^{r+s}_{Z}(X, F).
\]
由定理~\ref{thm:16.7}，知道除非 $s = 2c$，$E^{r, s}_2 = 0$.由此推出 $E^{r, s}_1 = E^{r, s}_2$，从而
\[
H^{r}_{Z}(X, F) \cong H^{r-2c}(Z, R^{2c}i^{\dagger}F) \cong H^{r-2c}(Z, F(-c)).
\]

\textbf{为什么定理~\ref{thm:16.7} 成立} 一旦定义了映射，问题对 $X$ 上的平展拓扑就是局部的.下一个引理表明任何余维数 $c$ 的光滑对都(对平展拓扑)局部同构于标准光滑对 $(\mathbb{A}^{m-c}, \mathbb{A}^{m})$.

\begin{lemma}\label{lem:16.8}
设 $Z \subset X$ 是余维数 $c$ 的光滑对，$P \in Z$.存在 $P$ 的开邻域 $V$ 与平展映射 $V \to \mathbb{A}^{m}$($m = \dim X$)，其到 $V \cap Z$ 的限制是到 $\mathbb{A}^{m-c}$ 的平展映射.
\end{lemma}

\begin{proof}
按假设，$\operatorname{Tgt}_{P}(Z)$ 是 $\operatorname{Tgt}_{P}(X)$ 的余维数 $c$ 的子空间.选取定义在 $P$($X$ 中)的开邻域 $V$ 上的正则函数 $f_1, \ldots, f_m$，使得 $df_1, \ldots, df_{m-c}$ 构成 $\operatorname{Tgt}_{P}(Z)$ 的对偶空间的一组基，且 $df_1, \ldots, df_m$ 构成 $\operatorname{Tgt}_{P}(X)$ 的对偶空间的一组基.考虑映射
\[
\alpha \colon V \to \mathbb{A}^{m}, \qquad Q \mapsto (f_1(Q), \ldots, f_m(Q)).
\]
映射 $d\alpha \colon \operatorname{Tgt}_{P}(X) \to \operatorname{Tgt}_{\alpha(P)}(\mathbb{A}^{m})$ 是 $v \mapsto (df_1(v), \ldots, df_m(v))$，它是同构.因此 $\alpha$ 在 $P$ 处平展.类似地，它到 $Z \to \mathbb{A}^{m-c}$ 的限制在 $P$(看作 $Z$ 的点)处平展.
\end{proof}

定理~\ref{thm:16.7} 的证明作归纳，从 $m = 1 = c$ 的情形出发，后者在 (\ref{prop:14.3}) 中证明.

\subsection*{推广}

回忆概形 $X$ 称为正则的，如果对所有 $x \in X$，局部环 $\mathcal{O}_{X, x}$ 正则.有限型态射 $\varphi \colon X \to S$ 称为光滑的，如果
\begin{enumerate}
  \item 它平坦，且
  \item 对每个代数闭几何点 $\bar{s} \to S$，纤维 $X_{\bar{s}}$ 正则.
\end{enumerate}
第二个条件意味着，对任何 $s \in S$，纤维 $\varphi^{-1}(s)$ 是非奇异簇.当 $X \to S$ 光滑时，我们说 $X$ 在 $S$ 上光滑.

对域 $k$ 上的簇 $X$，我们说 $X$ 非奇异，当 $X \to \operatorname{Spec} k$ 光滑时.于是，非奇异簇正则，而代数闭域上的正则簇光滑.非完全域上的簇可以正则而不非奇异---例如，若 $t \notin k$ 且 $p = \operatorname{char}(k)$，则 $Y^{2} = X^{p} - t$ 是这样的簇.

\begin{example}\label{ex:16.9}
设 $X_1$ 是 $\mathbb{Q}_p$ 上的非奇异完备曲线.则存在紧合平坦态射 $X \to \operatorname{Spec} \mathbb{Z}_p$，其一般纤维是 $X_1 \to \operatorname{Spec} \mathbb{Q}_p$.事实上，总是存在正则的这样的 $X$，但并非总是存在在 $\operatorname{Spec} \mathbb{Z}_p$ 上光滑的 $X$.对椭圆曲线，这在我们的椭圆曲线讲义 \S 9 中讨论.
\end{example}

如下猜想一般仍开放.

\begin{conjecture}\label{conj:16.10}
设 $X$ 是正则概形，$i \colon Z \hookrightarrow X$ 是 $X$ 的正则闭子概形，使得 $Z$ 在 $X$ 的每一点处都有余维数 $c$.则对所有满足 $n\mathcal{O}_X = \mathcal{O}_X$ 的整数 $n$，$R^{2c}i^{\dagger}\Lambda \cong \Lambda(-c)$，且对 $r \neq 2c$，$R^{r}i^{\dagger}\Lambda = 0$.
\end{conjecture}

当 $Z$ 与 $X$ 都在基概形 $S$ 上光滑时，该猜想在 SGA 4 中证明.在 $X$ 与 $n$ 的一些假设下，Thomason 能够从 Quillen 局部化定理(它是 $K$ 群的纯性论断)推出该猜想(Thomason, R.W., \emph{Absolute Cohomological Purity}, Bull. Soc. Math. France 112 (1984), 397--406).

\begin{proposition}\label{prop:16.11}
设 $(Z, X)$ 是纯性猜想 (\ref{conj:16.10}) 成立的一对余维数 $c$ 的对象.设 $U = X \setminus Z$，$i, j$ 如常表示包含映射.则
\[
R^{r}j_{*}\Lambda \cong \begin{cases}
\Lambda, & r = 0;\\
i_{*}\Lambda(-c), & r = 2c - 1;\\
0, & \text{其他情形}.
\end{cases}
\]
\end{proposition}

\begin{proof}
空间对 $(X, U)$ 的正合序列是
\[
\cdots \to H^{r}_{Z}(X, \Lambda) \to H^{r}(X, \Lambda) \to H^{r}(U, \Lambda) \to \cdots.
\]
对任何平展映射 $\varphi' \colon X' \to X$ 得到类似序列，只把 $U$ 与 $Z$ 换成 $U' = \varphi'^{-1}(U)$ 与 $Z' = \varphi'^{-1}(Z)$.当 $X'$ 变动时，这变成预层的正合序列，在应用``层化''函子 $a$ 后仍正合.现在与 $X' \mapsto H^{r}(X', \Lambda)$ 相伴的层对 $r > 0$ 为零，对 $r = 0$ 为 $\Lambda$，与 $X' \mapsto H^{r}(U', \Lambda)$ 相伴的层是 $R^{r}j_{*}\Lambda$.最后，与 $X' \mapsto H^{r}_{Z'}(X', \Lambda)$ 相伴的层是 $R^{r}i^{\dagger}\Lambda$，对 $r = 2c$ 它是 $\Lambda(-c)_Z$，否则为零.
\end{proof}

\begin{quote}
\textbf{注记} 关于纯性的更一般结果，见 Jo\"el Riou. \emph{Puret\'e (d'apr\`es Ofer Gabber)}. Notes des expos\'es des 16 juin 2006 et 13 avril 2007 au groupe de travail sur les r\'esultats r\'ecents de Gabber \`a l'\'Ecole Polytechnique, 2007 \href{http://www.math.u-psud.fr/~riou/doc/gysin.pdf}{此处}.
\end{quote}
